\section{Implementation Details and Detailed Experimental Results}\label{app:implementation}

\subsection{Probabilistic State Compression}\label{app:compression}

Our implementation reduces the state space of the lifted MDP by repeatedly applying \Cref{theo:compression} to states which where lifted from (inner) nodes with probabilistic variables. 
This theorem formally specifies the Compression procedure, and all target states always fulfill the requirements listed in \Cref{theo:compression}.
The soundness of each compression step follows from the theorem below, and we could apply the theorem exhaustively, but we enforce a maximum number of outgoing transitions for each state in the MDP, which guarantees that the Compression of the MDP as a whole can be completed in linear time relative to the size of the MDP.
In our implementation, that limit is set at $40$ transitions. 

\begin{theorem}[Probabilistic State Compression]\label{theo:compression}
	Let $M = (S, s_{init}, \mathit{Act}, P, \mathit{AP}, L)$ be an MDP and let $s_p \in S$ be a state with only one enabled action $\alpha_p$, as well as $s_p \neq s_{init}$, $P(s_p, \alpha_p, s_p) = 0$ and $L(s_p) = \emptyset$. 
	Now let $M' = (S/ \{s_p\}, s_{init}, \mathit{Act}, P', \mathit{AP}, L)$ be the compressed MDP where we define
	$$ P'(s, \alpha, s' ) \coloneqq P(s, \alpha, s') + \left(P(s, \alpha, s_p) \cdot P(s_p, \alpha_p, s')\right) $$
	We then have for all labels $v \in \mathit{AP}$:
	$$Pr^{max}_{M}(\lozenge v) = Pr^{max}_{M'}(\lozenge v)$$
\end{theorem}
\begin{proof}
	We once again consider the maximum reachability as characterized through Bellman equations \cite{forejt2011automated}. 
	To this end, we first need to define $S^0_M \coloneqq \{s \in S \,|\, Pr^{max}_M (\lozenge v) = 0\}$ to characterize the set of states from which no state labeled $v$ can be reached.
	Note that $S^0$ can be determined using graph based algorithms such as e.g. breadth first search, as it suffices to know if transitions are possible, while their probabilities can be ignored\cite{forejt2011automated}.
	Afterwards, the Bellman equations themselves are defined as follows, where $x(s)$ describes the probability of reaching a state labeled $v$ from the state $s$.
	$$ x(s) = \begin{cases}
		1 &, v \in L(s)\\
		0 &, s \in S^0_M\\
		\max\limits_{\alpha \in \mathit{Act(s)}} \sum\limits_{s' \in S} P(s, \alpha, s') \cdot x(s') &,\text{otherwise}
	\end{cases}$$
	$Pr^{max}_M (\lozenge v)$ is then the value for the initial state by the \textit{least} solution to the system of equations described by the Bellman Operator.
	We will write $x(s)$ for the Bellman equation for $M$ and $x'(s)$ for the Bellman equation of the compressed MDP $M'$.
	
	We will now show that for all states $s \in S/ \{s_p\}$ we have $x(s) = x'(s)$.
	This is trivial for all states $s$ with $v \in L(s)$ and follows for all states $s\in S^0$ from $S^0_M = S^0_{M'}$.
	We now consider the remaining case. We first show:
	\begin{align*}
		&\sum\limits_{s' \in S} P(s, \alpha, s') \cdot x(s') \\
		&= P(s, \alpha, s_p) \cdot x(s_p) + \sum\limits_{s' \in S/\{s_p\}} P(s, \alpha, s') \cdot x(s') \\
		&=_1 P(s, \alpha, s_p) \cdot \sum\limits_{s' \in S} P(s_p, \alpha_p, s') \cdot x(s') + \sum\limits_{s' \in S/\{s_p\}} P(s, \alpha, s') \cdot x(s') \\
		&=_2 P(s, \alpha, s_p) \cdot \sum\limits_{s' \in S/\{s_p\}} P(s_p, \alpha_p, s') \cdot x(s') + \sum\limits_{s' \in S/\{s_p\}} P(s, \alpha, s') \cdot x(s') \\
		&= \sum\limits_{s' \in S/\{s_p\}} P(s, \alpha, s_p) \cdot P(s_p, \alpha_p, s') \cdot x(s') + \sum\limits_{s' \in S/\{s_p\}} P(s, \alpha, s') \cdot x(s') \\
		&= \sum\limits_{s' \in S/\{s_p\}} \left(P(s, \alpha, s') + P(s, \alpha, s_p) \cdot P(s_p, \alpha_p, s')\right) \cdot x(s') \\
		&= \sum\limits_{s' \in S/\{s_p\}} P'(s, \alpha, s') \cdot x(s') \\
	\end{align*}
	In $=_1$ we use that $L(s_p) = \emptyset$, meaning that the Bellman equation for $s_p$ is described by the final rule.
	We also use that $s_p$ only has a single enabled action, which simplifies the Bellman equation to just a sum as the maximum over all actions is trivially found with $\alpha_p$.
	Note that $x(s_p) = \sum_{s' \in S} P(s_p, \alpha_p, s') \cdot x(s')$ also holds in the case of $s_p \in S^0_M$, as we then either have $s' \in S^0_M$ or $P(s_p, \alpha_p, s') = 0$ for all states $s' \in S$, meaning the sum as whole will be $0$.
	Afterwards, $=_2$ follows from $P(s_p, \alpha_p, s_p) = 0$.
	It now follows:
	\begin{align*}
		x(s) & = \max\limits_{\alpha \in \mathit{Act(s)}} \sum\limits_{s' \in S} P(s, \alpha, s') \cdot x(s') \\
		&= \max\limits_{\alpha \in \mathit{Act(s)}} \sum\limits_{s' \in S/\{s_p\}} P'(s, \alpha, s') \cdot x(s') \\
	\end{align*}
	Therefore, the system of equations as defined by $x(s)$ is the same as the system of equations as defined by $x'(s)$ up to renaming of $x(s)$. 
	Therefore, the least solution to one system of equations is also the least solution for the other system of equations, and we get $x(s) = x'(s)$ for all states $s \in S/\{s_p\}$.
	Finally, we have $s_p \neq s_{init}$ by assumption, and we can conclude $Pr^{max}_{M}(\lozenge v) = x(s_{init}) = x'(s_{init}) = Pr^{max}_{M'}(\lozenge v)$
\end{proof}



\subsection{Reduction to Boolean Inference}\label{app:boolean_reduction}
In this section we show that we can soundly reduce the semantics of arbitrarily typed programs to boolean programs.
The key idea is that we wrap the expression $e$ we are interested in into a let statement $\LET{x = e}{ x \leftrightarrow v}$, called the boolean reduction of $e$ for $v$, which essentially evaluates $e$ and then tests if the execution yielded the desired value $v$.
The idea is that it does not matter if we ask for the maximum probability for $e$ to return a particular value, or for the maximum probability for the boolean reduction of $e$ for $v$ to return $\TRUE$.
To this end we define a simple equivalence or equality test on values with the following semantics:
$$ \llbracket v_1 \leftrightarrow v_2 \rrbracket^T \Coloneqq \begin{cases}
	\TRUE| 1\rangle &, v_1 = v_2 \\
	\FALSE| 1\rangle &, v_1 \neq v_2 \\
\end{cases}$$
With this definition in hand we present \Cref{theo:boolean_reduction}, which relates the (unnormalized) distributions in the semantics of expressions and their boolean reductions to each other.
Afterwards, \Cref{theo:boolean_reduction_norm} then considers normalized distributions and proves the soundness of this inference approach.
\begin{lemma}[Unnormalized Boolean Reduction]\label{theo:boolean_reduction}
	Let $e$ be an \NDICE expression, and let $T$ be a function table.
	Then for all values $v \in V$ we have:
	\begin{enumerate}
		\item $\forall d \in \llbracket e \rrbracket^T \; \exists d' \in \llbracket \LET{x = e}{ x \leftrightarrow v}\rrbracket^T: d(v) = d'(\TRUE) \land \sum\limits_{w \neq v} d(w) = d'(F)$
		\item $\forall d' \in \llbracket \LET{x = e}{ x \leftrightarrow v} \rrbracket^T \; \exists d \in \llbracket e \rrbracket^T: d(v) = d'(\TRUE) \land \sum\limits_{w \neq v} d(w) = d'(F)$
	\end{enumerate}
\end{lemma}
\begin{proof}
	Let $v \in V$ be some value.
	\begin{enumerate}
		\item Let $d \in \llbracket e \rrbracket^T$ be some distribution.
		Then, by the semantics of the let-statement in \Cref{fig:semantics} we have
		$$d' \in \llbracket \LET{x = e}{ x \leftrightarrow v} \rrbracket^T \Leftrightarrow \exists d_1 \in \llbracket e \rrbracket^T: \forall w, d_1(w) > 0: \,\exists d_{w} \in \llbracket x\leftrightarrow v [x \mapsto w] \rrbracket^T: $$
		$$ \forall u \in V : d'(u) = \sum_{w,\, d_1(w) > 0} d_1(w) \cdot d_{w}(u)$$
		Note that $\llbracket x\leftrightarrow v[x \mapsto w] \rrbracket^T = \llbracket w\leftrightarrow v \rrbracket^T$ is always a singleton set. 
		We can therefore pick $d_1$ to be $d$, and pick $d_w$ as the only element of $\llbracket w\leftrightarrow v \rrbracket^T$ to define the distribution $d' \in \llbracket \LET{x = e}{ x \leftrightarrow v} \rrbracket^T$.
		We first consider both $d(v)$ and $d'(\TRUE)$:
		\begin{align*}
			d'(\TRUE) &= \sum_{w,\, d_1(w) > 0} d_1(w) \cdot d_{w}(\TRUE) \\
			&= \sum_{w,\, d(w) > 0} d(w) \cdot d_{w}(\TRUE) \\
			&=_1 d(v) \cdot d_{v}(\TRUE) \\
			&= d(v)
		\end{align*}
		Note that the final two equations hold as $d_w(\TRUE) = 0$ for all $w \neq v$, and $d_w(\TRUE) = 1$ for $w = v$ by definition.
		Finally, notice that $=_1$ holds both if we have have $d(v) > 0$ or not. 
		If we have $d(v) > 0$, then $=_1$ holds as argued above because $d(v) \cdot d_v(\TRUE)$ is a proper term in the sum and is also the only non-zero summand.
		If we have $d(v) = 0$, then $=_1$ holds as both sides of the equation are $0$, in particular the sum is zero because all the $d_w(\TRUE)$ are $0$.
		
		The second part of the conjunction can be proven similarly.
		However, this time we have  $d_w(\FALSE) = 0$ for all $w = v$, and $d_w(\FALSE) = 1$ for $w \neq v$, allowing us to show the following equations:
		\begin{align*}
			d'(\FALSE) &= \sum_{w,\, d_1(w) > 0} d_1(w) \cdot d_{w}(\FALSE) \\
			&= \sum_{w,\, d(w) > 0 \land w \neq v} d(w) \cdot d_{w}(\FALSE) \\
			&= \sum_{ w \neq v} d(w) \cdot d_{w}(\FALSE) \\
			&= \sum_{ w \neq v} d(w)
		\end{align*}
		
		
		\item Let $d' \in \llbracket \LET{x = e}{ x \leftrightarrow v} \rrbracket^T$.
		Therefore, by the semantics of the let-statement as defined in \Cref{fig:semantics} there exists a distribution 
		$d \in \llbracket e \rrbracket^T$, and for all values $w$ with $d(w) > 0$ there exists a distribution $d_{w} \in \llbracket x\leftrightarrow v [x \mapsto w] \rrbracket^T$ so that 
		$$ \forall u \in V : d'(u) = \sum_{w,\, d(w) > 0} d(w) \cdot d_{w}(u)$$
		We now consider $d'(\TRUE)$ and $d(v)$.
		Once again, we note that $\llbracket x\leftrightarrow v[x \mapsto w] \rrbracket^T = \llbracket w\leftrightarrow v \rrbracket^T$ is always a singleton set, meaning we can simplify this sum in accordance with the semantics of $w\leftrightarrow v$ as follows:
		$$ d'(\TRUE) = \sum_{w,\, d(w) > 0} d_1(w) \cdot d_{w}(\TRUE) =_2 d(v) \cdot d_v(\TRUE) = d(v) $$
		These equations are justified very similarly to the case above, where we note that $d_w(\TRUE) = 0$ for all $w \neq v$, and $d_w(\TRUE) = 1$ for $w = v$ by definition.
		Furthermore, the second part of the conjunction can be proven analogously to $(1)$ as well.
	\end{enumerate} 
\end{proof}


\begin{theorem}[Normalized Boolean Reduction]\label{theo:boolean_reduction_norm}
	Let $e$ be an \NDICE expression, and let $T$ be a function table.
	Then for all values $v \in V$ we have:
	\begin{enumerate}
		\item $\forall d \in \llbracket e \rrbracket_D^T \; \exists d' \in \llbracket \LET{x = e}{ x \leftrightarrow v}\rrbracket_D^T: d(v) = d'(\TRUE) \land \sum\limits_{w \neq v} d(w) = d'(F)$
		\item $\forall d' \in \llbracket \LET{x = e}{ x \leftrightarrow v} \rrbracket_D^T \; \exists d \in \llbracket e \rrbracket_D^T: d(v) = d'(\TRUE) \land \sum\limits_{w \neq v} d(w) = d'(F)$
	\end{enumerate}
\end{theorem}
\begin{proof}
	The theorem follows directly from \Cref{theo:boolean_reduction}, which correlates the unnormalized distributions in the semantics of both expressions.
	Afterwards, the normalizing constants for corresponding distributions can also be shown to be equal, which directly proves the required relation between the normalized semantics of both expressions
\end{proof}

\subsection{Extended Experimental Results}\label{app:extended_results}

\setboolean{useComma}{true}

This section extends the empirical results from \Cref{sec:implementation} with the standard deviation over the reported runtimes, all standard deviations are measured over at least 5 evaluations of the respective benchmark.
We do not report standard deviations for the sizes of the BDDs, ADDs and MDPs again as their numbers are constant between all runs of a benchmark.
Below, \Cref{tab:Ndet_extended} extends \Cref{tab:Ndet}.

\begin{table}[H]
	\caption{Performance of \NDICE on the benchmarks. 
		The \DICE, \NDICE and Storm columns report the total runtime of the respective algorithm, \NDICE is run with state compression and split into \textit{compilation} and \textit{model checking}.
		$\NDICE_u$ reports the total inference time without state compression.
		$\NDICE_u$ reports the total inference time without state compression.
		The reported time is the mean plus or minus a single standard deviation over 5 runs.
		All times are in ms, with a timeout (TO) after 20 minutes.
		}
	\label{tab:Ndet_extended}
	\begin{adjustbox}{max width=\textwidth}
	\begin{tabu}{c | c | ccc | c | c}
		\toprule
		Benchmark & \DICE &  $\mathit{compilation}$ & $\textit{checking}$ & \NDICE & $\NDICE_u$ & Storm \\
		\midrule
		Runway-7 & - & $8\pm 0.4$ & {$\approx 0 \pm 0$} & {$8 \pm 0.4$} & {$8 \pm 0.4$} & {$18 \pm 2.8$} \\
		Runway-15 & - & $93 \pm 2.3$ & {$1 \pm $} & {$94 \pm 2$} & {$94 \pm 2 $} & {$131 \pm 2.7$} \\
		Runway-30 & - & $531 \pm 7$ & {$7 \pm 0.1 $} & {$538 \pm 7.1 $} & {$538 \pm 7.5$ } & {$1\, 620 \pm 14$} \\
		Runway-40 & - & $999 \pm 8$ & {$11 \pm 0.5 $} & {$1\,010 \pm 8.2$} & {$1\,011 \pm  8.8$} & {$5\,212 \pm 61$} \\
		Runway-45 & - & $1\,514 \pm 19$ & {$15 \pm 0.4$} & {$1\,529 \pm 18$} & {$1\,530 \pm 20 $} & {$8\,673 \pm 66$} \\
		
		\midrule
		CouponProb-10 & 
		$721 \pm 11$ &
		$1\,379 \pm 46$ & {$\approx 0 \pm 0$} & {$1\,379 \pm 46$} & 
		{$1\,379 \pm 46$} & $5\,894 \pm 106$ \\
		CouponProb-11 & 
		$1\,958 \pm 33$ & 
		$3\,904 \pm 62$ & {$\approx 0 \pm 0$} & {$3\,904 \pm 62$} & 
		{$3\,904 \pm 62$} & $18\,839 \pm 206$ \\
		CouponProb-12 & 
		$5\,798 \pm 161$ & 
		$10\,484 \pm 216$ & {$\approx 0 \pm 0$} & {$10\,484 \pm 216$} & 
		{$10\,484 \pm 217$} & $56\,434 \pm 1\,115$ \\
		CouponProb-13 & 
		$20\,691 \pm 106$ & 
		$32\,153 \pm 546$ & {$\approx 0 \pm 0$} & {$32\,153 \pm 546$} & 
		{$32\,153 \pm 546$} & $171\,358 \pm 442$ \\
		CouponProb-14 & 
		$63\,418 \pm 280$ & 
		$96\,068 \pm 647$ & {$\approx 0 \pm 0$} & {$96\,068 \pm 647$} & 
		{$96\,068 \pm 647$} & $641\,614 \pm 231\,908$ \\
		
		\midrule
		CouponNdet-6 & 
				- & $105 \pm 1$ &  {$14 \pm 0.7$} & {$119 \pm 1.3$} & {$120 \pm 1.8$} & $162 \pm 7.6$ \\
		CouponNdet-7 & 
				- & $276 \pm 3$ & {$41 \pm 1.8$} & {$317 \pm 2.5$} & {$317 \pm 3.4$} & $607 \pm 25$ \\
		CouponNdet-8 & 
				- & $731 \pm 19$ & {$96 \pm 2$} & {$827 \pm 20$} & {$824 \pm 19$} & $2\,208 \pm 99$ \\
		CouponNdet-9 & 
				- & $2\,264 \pm 49$ & {$389 \pm 33$} & {$2\,653 \pm 66$} & {$2\,662 \pm 64$} & $8\,110 \pm 23$ \\
		CouponNdet-10 & 
				- & $6\,105 \pm 107$ & {$1\,257 \pm 81$} & {$7\,362 \pm 114$} & {$7\,293 \pm 117$} & $25\,862 \pm 388$ \\
		 
		\midrule
		Network-2000 & - & $2\,218 \pm 16$ & {$24 \pm 0.9$} & {$2\,242 \pm 17$} & {$2\,242 \pm 17$} & $25 \pm 0.4$ \\
		Network-4000 & - & $9\,530 \pm 46$ & {$55 \pm 3$} & {$9\,585 \pm 45$} & {$9\,585 \pm 48$} & $39 \pm 0.4$ \\
		Network-6000 & - & $22\,270 \pm 98$ & {$84 \pm 1.3$} & {$22\,354 \pm 98$} &{$22\,358 \pm 102$} & $54 \pm 1.1$ \\
		Network-8000 & - & $40\,763 \pm 98$ & {$122 \pm 6$} & {$40\,885 \pm 97$} & {$40\,885 \pm 98$} & $68 \pm 1.6$ \\
		Network-10000 & - & $67\,081 \pm 207$ & {$153 \pm 1.7$} & {$67\,234 \pm 255$} & {$67\,243 \pm 214$} & $84 \pm 1.4$ \\
		
		\midrule
		Rabin-50 & - & $324\pm 1.1$ & {$\approx 0 \pm 0$} & {$ 324 \pm 1.1$} & {$ 324 \pm  1.1$} & $17 \pm 4.5$ \\
		Rabin-100 & - & $1\,962\pm 49$ & {$1 \pm 0$} & {$ 1\,963 \pm  49$} & {$ 1\,963 \pm  49$} & $23 \pm 0.7$ \\
		Rabin-250 & - & $23\,355\pm 141$ & {$2 \pm 0$} & {$ 23\,357 \pm  141$} & {$ 23\,357 \pm  142$} & $82 \pm 3.5$ \\
		Rabin-500 & - & $233\,494\pm 1\,759$ & {$6 \pm 0.5$} & {$233\,500 \pm 1\,759$} & {$233\,501 \pm 1\,760$} & $327 \pm 6.9$ \\
		
		\midrule 
		Survey & $1 \pm 0.2$ & $1 \pm 0$ & $\approx 0 \pm 0$ & $1 \pm 0$ & $1 \pm 0$ & $14 \pm 2$  \\
		Hepar2 & $99 \pm 1$ & $309 \pm 3$ & $\approx 0 \pm 0$ & $309 \pm 3$ & $310 \pm 3$ & $23\,094 \pm 146$ \\
		Insurance & $580 \pm 4$ & $2\,227 \pm 13$ & $\approx 0 \pm 0$ & $2\,227 \pm 13$ & $2\,232 \pm 14$ & $9\,460 \pm 146$ \\
		Water & $9\,241 \pm 47$ & $9\,288 \pm 37$ & $\approx 0 \pm 0$ & $9\,288 \pm 37$  & $9\,347 \pm 37$   & $12\,545 \pm 46$ \\
		Alarm & $160\,104 \pm 367$ & $161\,385 \pm 533$ & $\approx 0 \pm 0$ & $161\,385 \pm 533$ & $161\,385 \pm 532$ & $5\,335 \pm 9$ \\
		
		\midrule 
		Survey-Ndet & - & $1 \pm 0.1$ & $\approx 0 \pm 0$ & $1 \pm 0.1$ & $1 \pm 0.1$ & $14 \pm 0.5$ \\
		Hepar2-Ndet & - & $280 \pm 2.6$ & $14 \pm 0.4$ & $294 \pm 2.9$ & $295 \pm 3.1$ & $23\,909 \pm 131$ \\
		Insurance-Ndet & - & $2\,217 \pm 12$ & $\approx 0 \pm 0$ & $2\,217 \pm 12$ & $2\,220 \pm 12$ & $10\,345 \pm 137$ \\
		Water-Ndet & - & $9\,296 \pm 97$ & $\approx 0 \pm 0$ & $9\,296 \pm 97$ & $9\,298 \pm 98$ & $13\,153 \pm 60$ \\
		Alarm-Ndet & - & $161\,786 \pm 1\,166$ & $47 \pm 2$ & $161\,633 \pm 1\,167$ & $161\,636 \pm 1\,166$ & $5\,219 \pm 17$ \\
		
		\midrule
		3sat-0 & $19\,511 \pm 72$ & $19\,535 \pm 73$ & {$\approx 0 \pm 0$} & {$19\,535 \pm 73$} & {$19\,535 \pm 73$} & TO \\
		3sat-25 & - & $19\,528 \pm 29$ & {$43 \pm 1$} & {$19\,571 \pm 29$} & {$19\,571 \pm 31$} & TO \\
		3sat-50 & - & $19\,544 \pm 58$ & {$69 \pm 0.4$} & {$19\,613 \pm 29$} & {$19\,612 \pm 30$} & TO \\
		3sat-75 & - & $19\,514 \pm 63$ & {$72 \pm 1.5$} & {$19\,586 \pm 64$} & {$19\,584 \pm 130$} & TO \\
		3sat-100 & - & $19\,567 \pm 128$ & {$107 \pm 3$} & {$19\,674 \pm 128$} & {$19\,676 \pm 132$} & TO \\
		\bottomrule 
	\end{tabu}

\end{adjustbox}

\end{table}
